\documentclass[script,con]{debate}
\motion{对宠物的爱，爱的更是它 / 爱的更是我}
\stance{反方}{爱的更是我}
\meta{赛制}{中国科学技术大学新生辩论赛（四对四，正方先立论，反方先结辩）}
\meta{口径来源}{备赛文档 第十节 10.2 反方口径表。本文稿所有定义、判准、论点、数据均取自该表}
\begin{document}
\debatetitle
\begin{thesisbox}[全场一句话立论]问卷上愿意负债的是九成，账单真来时让它等过的近半；病老之时它一句话说不了，多疼一段还是早走一步，跟着我的不舍和我的负担走。冲突时让路的是它，决定的是我，所以爱的更是我。\end{thesisbox}
\begin{thesisbox}[全队第一纪律]每个环节至少说一次“我方不否认你爱它”。这不是客套，这是让评委愿意继续听下去的门票。\end{thesisbox}

\begin{speech}{1. 反一开篇立论稿}{正文 773 字 / 180 秒}
谢谢主席。

刚才的质询里，正方一辩已经承认了两件事：九成愿意负债是一道意愿题；行为题里，让自己的人和让它等的人，两边都不过半。请评委先把这两句记下来。

我先说清楚一件事，免得后面被误会。我方不认为养宠的人自私，更不认为你不爱它。我方今天要说的从来不是“你不爱”，而是这份爱的重心落在哪一边。

定义上我方和正方没有分歧。爱是对一个对象的深挚感情，以及这份感情带来的关切和行动。判准我方也接受正方那一条，冲突时优先谁。我方只加一句：以行为为准，以多数为准。谁做了什么，比谁说了什么，更能看出让路的方向。

我方第一点：说的是为它，做的跟着我的钱走。

正方引的那份两千人报告，我方也读了。同一份报告，正方念的是意愿题，九成一的人说愿意为救它负债。我方念的是行为题，往下翻一页：四成六的人在过去一年因为费用，延迟或者跳过了它的就医；四成四的人上一次带它看病，差一点因为钱放弃。心理学管这个叫意图和行为之间的鸿沟。说的时候我们都是好主人，账单真的来了，先等的是它的病。它等不等，看的不是它有多疼，看的是我的余额。我方承认，同一份报告里也有三成七的人削了自己的口粮；行为题两边都不过半，正方要的“多数”不在这里。

我方第二点：病老之时，舍不得就让它多受一段。

正方说，回报归零的时候投入最高。请评委看清楚那个时候谁在做决定。它一句话说不了。一项十七个国家两百二十八位宠物主的研究里，安乐死决策的主导因素是什么？是主人的情感依恋，近七成。真正拿它的生活质量指标来判断的，只有百分之七。正方引的那位学者，同一个人二零二二年发现，同样的疼痛，主人的负担越高，越早考虑放手。倒过来也一样：八百八十九位北美兽医里近八成被要求过继续做他们自己都认为没用的治疗。舍不得，它就多受一段。你想想，同一只狗，同样的疼，落在一个撑得住的主人和一个撑不住的主人身上，结局不一样。变的不是它的疼，变的是我。

所以我方全场要请对方回答一个问题：它多疼一段还是早走一步，你方说跟着它的指标走，那用指标的人在哪里？

我方承认这份感情真实。我方只是想说，它最深的那一层，跟着的还是我们自己。

谢谢。
\end{speech}
\begin{contingency}
\ifthen{正一没有承认“行为题两边不过半”}{开头改成：“刚才正方一辩没有否认，九成愿意负债是一道意愿题。”其余不动。}
\ifthen{正方问“有没有任何一种爱更指向对方”}{正面答“有”，母亲为孩子放弃职业规划；然后要条件，逐条搬到养宠里比。绝不能答“一种都没有”。}
\ifthen{正方定义“爱”为“为对方之故”}{30 秒口径：这是定义杀，请对方回答，被告知狗超重仍每周喂三次零食的主人算不算爱狗。}
\end{contingency}

\begin{stage}{2. 反一被质询防守预案}{正四质询，120 秒双边计时}
\textbf{原则}：凡是与常识冲突的都不否认，凡是道德指控都当场撇清。全场绝对不能说的三句：“养宠的人自私”“养宠的人在骗自己”“安乐死是它让路、继续治也是它让路”。
\begin{qa}
\qarow{母亲对孩子的爱，更是谁？}{更是孩子。我方承认存在更指向对方的爱。}
\qarow{那饲主做什么才算爱的更是它？}{三样：有得选时让自己扛；按它的指标而不是我的状况定终点；被告知后按它的需要改。}
\qarow{按它的指标放手，算不算为它？}{算。我方说的是，多数人的开关不是它的指标，用指标的只有百分之七。}
\qarow{三成七削口粮、三成七负债，是不是行为题？}{是。行为题两边都不过半，所以正方“多数人让自己”也没有证据。}
\qarow{急诊会拖的一成六比四成六少，对吗？}{那是假设题，和九成一同一题型。正方自己说假设题不证多数。}
\qarow{负担中介的起点是不是它的疼？}{是。可如果决定真跟着它走，负担这个中介就不该显著。同样的疼，不同的我，不同的决定。}
\qarow{安乐死理由八成七是它的疼，对吗？}{对，那是事后填的理由；同一张问卷里主导因素填的是依恋。两栏都是自报。}
\qarow{它病了，主人的负担变不变？}{变。我方不否认我的状况随它的病变；我方说决定跟着的是它的病落在我身上的量。}
\qarow{你方是不是在说人都自私？}{不是。我方全场没有一句道德判断，我方不否认你爱它。}
\qarow{你方的“拖”有饲主比例吗？}{没有，我方承认。正方“按指标及时放手”的比例也没有。有比例的是用指标的只有百分之七。}
\end{qa}
\end{stage}

\begin{stage}{3. 反四质询正一问题链}{120 秒，双边计时}
\textbf{总目标}：拿到“九成是意愿题”“行为题两边不过半”“被需要是我的需要”三个承认，反一开头要用。
\begin{chain}{链一：意愿题还是行为题}{拆掉正方论点二的“多数”}
\q{九成愿意负债，是意愿题还是行为题？}{答“意愿题”：下一问；答“行为题”：请念原文“愿意”。}
\q{同一份报告，四成六因成本延迟过就医，对吗？}{答“对”：下一问；答“不对”：请念原文。}
\q{行为题里让自己的人过半吗？}{答“不过半”：记住这句；答“过半”：请给数字。}
\close{正方承认九成是意愿题，行为题两边不过半，“多数人让自己”没有证据。}
\end{chain}
\begin{chain}{链二：被需要}{拆掉“回报归零”的前提}
\q{被需要，是不是你的需要？}{答“是”：下一问；答“不是”：共同定义第二项写着被需要。}
\q{它瘫在床上的时候，你是不是最被需要的时候？}{答“是”：那回报没有归零；答“不是”：那谁在需要你起床。}
\q{它治不治、走不走，是不是全由你决定？}{答“是”：可控感在峰值；答“不是”：请问还有谁。}
\close{被需要与可控感在病老期最高，“回报归零”在共同定义下不成立。}
\end{chain}
\begin{chain}{链三：开关在哪}{把正方的照护数据接到我方}
\q{负担越高越考虑安乐死，这条你方撤不撤？}{答“不撤”：下一问；答“撤”：那是你方引的作者。}
\q{同样的疼，负担不同，决定是不是不同？}{答“是”：开关在我；答“不是”：那中介为什么显著。}
\q{用结构化指标判断的只有百分之七，对吗？}{答“对”：收；答“不知道”：请评委记住。}
\cut{“好的，这个我理解为您没有否认。下一个问题。”}
\close{正方承认负担影响决定，用它的指标的只有百分之七。}
\end{chain}
\end{stage}

\begin{speech}{4. 反二驳论稿}{正文 505 字 / 120 秒}
谢谢主席。

正方立论最大的问题，是把“爱有代价”当成了“代价的方向是它”。

正方第一点说，回报归零的时候投入最高，证据是二百三十八位主人的对照研究：照护病宠的人负担更高、更抑郁。我方一点都不反对这个结论。可是请评委看清楚，这个量表测的是什么。它叫负担量表，题目问的是照护对你自己的生活造成了什么影响。它证明的是，照护让我很累。它没有一个变量叫投入，也没有一条时间线说投入不降反升。而且正方引的这位学者，同一个人二零二二年的研究说：负担越高的主人，越早考虑放手。请评委记住，负担不是撑着不放手的证据，负担是放手的预测变量。

正方第二点说，让路的是我，三成七削了口粮。我方承认，这是行为，这些人让了自己。可同一份报告里四成六让它等过，正方自己在质询里说了，行为题两边都不过半。那正方靠什么说多数人让自己？靠九成愿意负债。九成是意愿。老实说，一个人在问卷上说“我会为它负债”，说的是他想成为的那种主人；账单来了，四成六的人让它等了。说的是为它，做的跟着钱走。

正方还会说，没有人会为了爱自己，去选一段注定剧痛的关系。八成五的美国主人说，宠物是他们一年里最主要的快乐来源。十几年最大的快乐，换最后几个月的痛，你想想，这笔账哪个为自己打算的人不会签。

请评委回到判准。正方到现在为止，还没有回答我方的问题：它多疼一段还是早走一步，用它的指标判断的人在哪里。
\end{speech}
\begin{contingency}
\ifthen{正一稿没提“回报归零”}{把驳点一换成“被需要在病老期最高”，用质询链二的成果。}
\ifthen{正方先自己承认“行为题两边不过半”}{直接追问：那“多数人让自己”的证据是什么？只剩意愿题。}
\ifthen{正方指控我方犬儒}{立刻撇清：“我方承认你爱它，我方全场没有一句说养宠的人虚伪。”}
\end{contingency}

\begin{stage}{5. 反二对辩要点}{90 秒，单边计时}
\begin{points}{进攻点（每个 15–20 秒，说完就停）}
\item 九成是意愿题，行为题两边不过半。正方“多数人让自己”的证据在哪里？
\item 负担量表测的是照护对我自己生活的影响。同一位学者说负担高的先考虑放手。
\item 决策主导因素是依恋，近七成；用指标的只有百分之七。它的指标，谁读的？
\item 被需要在它瘫痪时最高，可控感也最高。回报归零了吗？
\item 八成五说它是最主要的快乐来源。明知结局仍然开始，是因为它值。
\end{points}
\begin{points}{备用点（回应对方进攻用）}
\item 我方承认你爱它，我方比的是分量，不是有无。
\item 母亲对孩子的爱更是孩子，我方承认存在更指向对方的爱。
\item 付不起也是相撞，先等的是它；让谁让路看我的余额。
\item 八成四是假设题，正方自己说假设题不证多数。
\item 我的状况随它的病变，我方不否认；决定跟着的是它的病落在我身上的量。
\item 理由栏写它疼，主导因素栏写依恋，两栏都是自报。
\item 兽医数据不是饲主比例，我方只用它说现象普遍。
\item 拖延的饲主比例双方都没有；有比例的是用指标的只有百分之七。
\end{points}
\begin{points}{对方最可能问的三个问题 → 回应}
\item “饲主做什么才算爱的更是它？” → 三样：有得选时让自己扛、按它的指标定终点、被告知后改。用指标的只有百分之七。
\item “负担中介的起点是它的疼吧？” → 是。如果决定跟着它走，中介就不该显著。
\item “你方是不是说人都自私？” → 不是，我方全场没有一句道德判断。
\end{points}
\end{stage}

\begin{stage}{6. 反三盘问问题链}{90 秒，单边计时}
\begin{chain}{链一：口径一致性检验}{让正二与正四在“负担影不影响它哪天走”上出现分歧}
\q{问正二：主人的负担，影不影响它哪天走？}{答“不影响”：那 2022 年那篇怎么解释；答“影响”：记住这句。}
\q{同一个问题问正四：您同意二辩的回答吗？}{答“同意”：口径锁死；答“不同意”：两人不一致。}
\q{那同样的疼，负担不同，决定是不是不同？}{答“是”：决定跟着我走；答“不是”：请解释中介为什么显著。}
\q{用它的指标判断终点的，是不是只有百分之七？}{答“是”：收；答“不知道”：请评委记住。}
\close{正方承认负担影响它哪天走，而用它的指标的只有百分之七。}
\end{chain}
\begin{chain}{链二：把“多数”逼出来}{证明正方在每个战场都拿不出过半的行为数据}
\q{问正四：行为题里让自己的人，过半吗？}{答“不过半”：下一问；答“过半”：请给数字。}
\q{病老期继续照护的人占多少，你方有比例吗？}{答“没有”：下一问；答“有”：请给。}
\q{问正二：按指标及时放手的饲主比例，你方有吗？}{答“没有”：收；答“有”：请给。}
\close{三个战场，正方一个过半的行为比例都没有。}
\end{chain}
\textbf{盘问目标}：第一，正方二辩四辩在“负担影不影响它哪天走”上要么口径不一，要么承认影响；第二，正方在三个战场都拿不出过半的行为比例。
\end{stage}

\begin{stage}{7. 反二、反四被盘问防守预案}{两人口径必须一致}
\begin{qaa}
\qaarow{母亲对孩子的爱更是谁？}{更是孩子。}{更是孩子，我方承认有这种爱。}
\qaarow{饲主做什么才算为它？}{三样：让自己扛、按指标定终点、被告知后改。}{同上三样。}
\qaarow{按它的指标放手算不算为它？}{算。用指标的只有百分之七。}{算。多数人的开关是依恋。}
\qaarow{三成七削口粮是不是行为题？}{是。两边都不过半。}{是。正方“多数”没有证据。}
\qaarow{负担中介的起点是不是它的疼？}{是。可中介显著说明决定跟着我的负担变。}{是。同样的疼，不同的我，不同的决定。}
\qaarow{它病了你的状况变不变？}{变。决定跟着的是它的病落在我身上的量。}{变。我方不否认代价，我方说方向。}
\qaarow{拖延有饲主比例吗？}{没有，我方承认。正方也没有“及时放手”的比例。}{同上。}
\qaarow{你方是不是说人都自私？}{不是，我方没有道德判断。}{不是，我方承认你爱它。}
\end{qaa}
两人赛前必须逐条对答一遍，答案可以换词，不能换意思。被逼到没准备过的问题时，统一说“我方立论的立场是”，不要现场发挥。
\end{stage}

\begin{speech}{8. 反三盘问小结稿}{正文 339 字 + 临场位 80 字 / 90 秒}
谢谢主席。

\reserve{此处留临场位，约 80 字：直接回应正三小结里最用力的一句。正三很可能会说我方的“不对称”被自己的数据否定。这里先接住：我方从没说我的状况不变，我方说的是决定跟着谁走。}
刚才的盘问里我问了两件事。第一件，主人的负担，影不影响它哪天走。请评委注意正方二辩和四辩的回答。说不影响，那正方自己引的那位学者二零二二年的研究就没法解释。说影响，那正方就承认了：同样的疼，负担不同，决定不同。开关在我。

第二件，我请正方在三个战场各给一个过半的行为比例。钱的战场，正方自己说两边不过半。病老的战场，继续照护的人占多少，正方没有。终点呢？按指标及时放手的人占多少，正方也没有，而且正方用来证明“多数人按它的指标来定终点”的那份研究，同一段摘要写着用生活质量指标判断的只有百分之七。

说白了，正方全场证明的是爱有代价。我方承认，我方不否认你爱它。可判准问的是撞上的时候让谁让路，代价说明爱是真的，说明不了方向。它等不等，看我的余额；它多疼一段还是早走一步，看我的不舍。

到这里，正方还是没有回答我方一辩问的那一句：用它的指标判断的人在哪里。

谢谢。
\end{speech}
\begin{contingency}
\ifthen{正三小结主打“不对称被自引数据否定”}{临场位正面接：我方从没说我的状况不变，我方说决定跟着谁走。}
\ifthen{正二正四口径一致且都说“影响”}{直接推进：“那请正方解释，负担影响它哪天走，开关怎么还在它身上。”}
\end{contingency}

\begin{stage}{9. 自由辩战场设计}{240 秒}
\begin{battleground}{战场一：意愿题还是行为题（前 90 秒，主战场）}
\begin{fields}
\field{目标}{让评委在同一份报告上看见“说的”和“做的”是两回事。}
\field{开场问题}{九成愿意负债，是意愿题还是行为题？}
\field{对方答“意愿题，但行为题有三成七” → 追问}{行为题两边不过半，那“多数人让自己”靠什么？}
\field{对方答“方向” → 追问}{方向来自意愿题，意愿题测的是你想成为哪种主人。}
\field{对方可能反攻 → 我方防守}{“急诊会拖的只有一成六。” → 假设题，和九成一同一题型，正方自己说假设题不证多数。}
\field{收束语}{说的是为它，做的跟着钱走。}
\end{fields}
\end{battleground}
\begin{battleground}{战场二：开关在哪（中 90 秒）}
\begin{fields}
\field{目标}{把正方的照护数据和终点数据转成我方的决定变量证据。}
\field{开场问题}{同样的疼，负担不同，决定是不是不同？}
\field{对方答“是” → 追问}{那开关在我的负担上，正方说的“开关在它的指标上”怎么成立？}
\field{对方答“渠道是我，起点是它” → 追问}{渠道的宽窄决定什么时候到，用指标的只有百分之七。}
\field{对方可能反攻 → 我方防守}{“理由八成七是它的疼。” → 理由是事后填的，同一张问卷主导因素填的是依恋，两栏都是自报。}
\field{收束语}{它多疼一段还是早走一步，跟着我的负担和不舍走。}
\end{fields}
\end{battleground}
\begin{battleground}{战场三：知道了改不改（最后 60 秒）}
\begin{fields}
\field{目标}{用日常画面收束，不新开战场。}
\field{开场问题}{它要少吃，你要它吃得开心，相撞了，让的是谁？}
\field{对方答“翻译错了不等于翻译给自己” → 追问}{被告知之后呢？被提示宠物超重的主人里六成一没有应对。}
\field{收束语}{它要散热，你要它穿得可爱；相撞了，让的是它的身体。}
\end{fields}
\end{battleground}
\begin{points}{必问（前两分钟内问完）}
\item 用它的指标判断终点的人在哪里？（一辩稿抛出的那个问题）
\item 九成愿意负债，是意愿题还是行为题？
\item 同样的疼，负担不同，决定是不是不同？
\item 它瘫在床上时，你是不是最被需要的时候？
\item 被告知它超重之后，改不改？
\end{points}
\begin{points}{必答（15 秒以内正面回答）}
\item “有没有任何一种爱更指向对方？” → 有，母亲为孩子放弃职业规划。请你方给条件，我们逐条比。
\item “饲主做什么才算为它？” → 三样：有得选时让自己扛、按它的指标定终点、被告知后改。用指标的只有百分之七。
\item “它病了你的状况变不变？” → 变。我方不否认代价，我方说决定跟着谁走。
\item “你方是不是说人都自私？” → 不是。我方承认你爱它，我方比的是分量。
\item “Zarit 量表是为照护失智父母的子女设计的，他们也更爱自己吗？” → 我方不用负担分判母爱，用它判决定跟着谁走；我方说的是分量，不是不爱。
\end{points}
\textbf{时间分配与站位}：前 90 秒战场一，抛出必问一、二；中 90 秒战场二，处理对方进攻；后 60 秒战场三，不开新战场，留 20 秒备用。一辩守定义与判准，二辩接驳论过的点，三辩接盘问成果，四辩收束与价值。
\end{stage}

\begin{speech}{10. 反四结辩稿}{正文 880 字 / 210 秒，封路式}
谢谢主席。

我先做一件事：把正方四辩待会儿要走的那条路，提前指出来。

正方四辩接下来大概率会给各位讲一个故事。一只很老的狗，失禁、瘫痪、痴呆，已经不认得主人了，主人一晚上起来六七次，后来辞了职。然后他会说，你看，在它什么都给不了的时候，这个人还在做，所以爱的是它。

这个故事我方相信是真的，讲的人也是真诚的。可是请评委听清楚，这里面有没被说出来的另一半。

那个每天喂药的人，为什么喂？因为停下来她受不了。喂药让她能对自己说一句“我尽力了”。这句话说给谁听？说给她自己。我方不是说这不高尚，我方是说，请评委不要把“她很痛苦地坚持”直接读成“她为的是它”。正方全场最硬的那组数据，二百三十八位主人的对照研究，测的正是这个：照护让我很累，很抑郁。它证明的是代价。代价说明爱是真的，说明不了方向。

现在回到判准。双方用的是同一条：冲突时优先谁。

钱的战场，正方自己在质询里承认了，行为题两边都不过半。正方靠什么说多数人让自己？靠九成愿意负债。九成是问卷上的意愿。说白了，那是我们想成为的那种主人。账单真的来了，四成六的人让它等了。

病老的战场，正方说回报归零。你想想，它瘫在床上、每晚叫醒你的时候，你是不是这辈子最被需要的时候？它治不治、走不走，是不是百分之百在你手里？回报没有归零，被需要和可控感在那一刻到了顶。

终点的战场，正方说多数人按它的痛来定。可判断它痛不痛的是谁？十七个国家两百二十八位主人里，决策的主导因素是主人的依恋，近七成；真正用它的生活质量指标判断的，只有百分之七。正方引的那位学者自己发现，同样的疼，负担高的主人先考虑放手。它多疼一段还是早走一步，跟着的是我的负担和我的不舍。

我方全场没有一句是在说养宠的人自私。我方每一个环节都说了，我方不否认你爱它。

我想请各位想一件很小的事。

你的手机相册里，大概有几百张它的照片。

它那边，一张你的都没有。

它不知道你叫什么，不知道你在哪里上班，不知道你昨天为什么回来得那么晚。你难过的时候抱着它，它其实不知道你在难过。

这不是它的错。它就是这样一种生命。

所以这份关系从一开始就不是对等的。我们投进去的想象、期待、委屈、盼望，它一样也接不住。它接不住，我们还是投。投给谁了？投给了那个被我们想象出来的它。

承认这一点，不是让我们爱得少一点。

恰恰相反。只有先承认那面镜子在，我们才有可能放下镜子，低下头，看一眼真正的它到底想要什么。

它可能只是想下楼多闻十分钟。

谢谢。
\end{speech}
\begin{contingency}
\ifthen{正方前场已经承认“行为题两边不过半”}{封路段加一句：“这句话正方自己说了三次，请评委数一数。”}
\ifthen{正四抢先在自由辩里讲了老狗的故事}{封路段前移到开头第一句，直接说“那个故事里没被说出来的另一半是”。}
\ifthen{正方指控我方犬儒}{在升华段前加一句：“我方全场没有说过一句‘别骗自己了’，我方说的是‘再看它一眼’。”}
\end{contingency}
\end{document}
